\chapter{Healthcare Costs}

\section*{Definition and Overview}
Healthcare costs are the expenses associated with medical care, including GP visits, specialists, tests, medicines, hospital treatment, and allied health services. In Australia, the cost of care is shared between the government, private health insurers, and individuals. Medicare covers many GP visits and public hospital treatment, while the Pharmaceutical Benefits Scheme (PBS) subsidises most prescription medicines. However, many people face out-of-pocket costs for specialist care, tests, dentistry, allied health, and some medicines. Understanding how healthcare is funded, checking costs before treatment, and knowing about subsidies can help people manage their expenses.

\section*{Epidemiology}
Most Australians use the health system regularly, and most are satisfied with access. Around 87\% of GP visits are bulk-billed, but bulk-billing rates for specialists are lower, and around one in five Australians report delaying or skipping care because of cost. Out-of-pocket costs average several hundred dollars per person per year, and around 3--5\% of Australians experience financial barriers to care. Public hospital care is free for public patients, and private health insurance covers some or all of the costs of private treatment. The cost of healthcare is a growing concern as out-of-pocket expenses rise.

\section*{Aetiology and Pathophysiology}
This chapter concerns health system finance rather than a disease. The main elements are:

\begin{itemize}
    \item \textbf{Medicare:} the national scheme that covers GP visits, many tests, and public hospital treatment
    \item \textbf{Pharmaceutical Benefits Scheme (PBS):} subsidises most prescription medicines
    \item \textbf{Private health insurance:} covers private hospital treatment and some out-of-hospital services
    \item \textbf{Bulk billing:} when the doctor accepts the Medicare benefit as full payment, leaving no out-of-pocket cost
    \item \textbf{Gap payments:} the difference between the doctor's fee and the Medicare benefit
    \item \textbf{Out-of-pocket costs:} payments for services not covered or only partly covered
    \item \textbf{Concessions:} lower costs for pensioners, concession card holders, and children
    \item \textbf{Public vs private:} public patients in public hospitals pay nothing, while private patients pay via insurance and gaps
\end{itemize}

\section*{Clinical Features}
Common situations where people encounter healthcare costs include:

\begin{itemize}
    \item GP visits, which are often bulk-billed
    \item Specialist appointments, which typically have out-of-pocket costs
    \item Pathology and imaging tests, which may attract gaps
    \item Prescription medicines on the PBS and over-the-counter medicines
    \item Dental care, which is largely out of pocket
    \item Allied health services, including physiotherapy and psychology
    \item Private hospital treatment and ambulance services
    \item Medicines and treatments not on the PBS
\end{itemize}

\section*{Differential Diagnosis}
Managing costs requires understanding what applies in each situation.

\begin{itemize}
    \item Bulk-billed versus gap-charged consultations
    \item PBS-subsidised versus non-subsidised medicines
    \item Public versus private hospital treatment
    \item In-hospital versus out-of-hospital insurance cover
    \item Medicare-covered services versus those not covered
    \item Concession eligibility and safety nets
    \item Urgent versus elective care costs
\end{itemize}

\section*{When to Seek Medical Care}
Cost should never be the reason to delay urgent care. For non-urgent care, people who are concerned about costs should discuss them with their doctor, ask about bulk billing, and explore subsidised services. Financial counselling and social work support are available in many hospitals.

\textbf{Call 000 (or local emergency services) for:}
\begin{itemize}
    \item Any medical emergency, regardless of cost or insurance status
    \item Chest pain, stroke symptoms, severe breathing difficulty, or major trauma
    \item Hospital emergency departments must treat emergencies regardless of ability to pay
\end{itemize}

\section*{Diagnosis and Investigations}
Planning for healthcare costs involves knowing your coverage.

\begin{itemize}
    \item \textbf{Check fees in advance:} ask about the cost of consultations, tests, and procedures before booking
    \item \textbf{Check Medicare benefits:} the Medicare schedule shows the benefit for each service
    \item \textbf{Check insurance cover:} understand what your private health insurance covers and any gaps
    \item \textbf{Ask about bulk billing:} many GPs and some specialists bulk bill
    \item \textbf{Check PBS status:} the pharmacist can advise the subsidised price of a medicine
    \item \textbf{Understand safety nets:} the Medicare and PBS safety nets reduce costs for high users
    \item \textbf{Ask about payment plans:} many providers offer them
\end{itemize}

\section*{Management}
Managing healthcare costs is about planning and asking the right questions.

\begin{itemize}
    \item \textbf{Use bulk-billed services:} choose bulk-billed GPs and, where possible, bulk-billed specialists
    \item \textbf{Use public hospital care:} public patients receive free hospital treatment
    \item \textbf{Check medicine prices:} ask about the PBS price, generic options, and cheaper alternatives
    \item \textbf{Consider private health insurance:} the right policy can cover private treatment, but check what it actually covers
    \item \textbf{Use concession cards:} pension and concession cards reduce costs
    \item \textbf{Register for safety nets:} the Medicare and PBS safety nets cap yearly costs for high users
    \item \textbf{Use community services:} community health centres, public dental services, and telehealth
    \item \textbf{Ask for a payment plan:} for large bills
    \item \textbf{Seek financial counselling:} available through hospitals and community services
\end{itemize}

\section*{Complications}
\begin{itemize}
    \item Delayed or skipped care because of cost, leading to worse health outcomes
    \item "Bill shock" from unexpected out-of-pocket costs
    \item Underinsurance, where policies do not cover what people expect
    \item Difficulty affording medicines, leading to non-adherence
    \item Financial stress and its effects on health
    \item Inconsistent access to care for people with lower incomes
    \item Confusion about what Medicare, the PBS, and insurance cover
\end{itemize}

\section*{Prognosis and Outcomes}
Australia's mixed funding system provides broad access to care, and most people receive the care they need. Understanding the system reduces out-of-pocket costs and improves access. People who ask about costs, use bulk billing and public services, and understand their insurance get better value from the system. For people facing genuine financial barriers, social workers, financial counsellors, and community health services provide support. Overall, the system delivers good outcomes, and cost should not be a barrier to essential care.

\section*{Prevention and Self-Care}
\begin{itemize}
    \item Ask about fees before appointments, tests, and procedures
    \item Choose bulk-billed GPs and services where possible
    \item Understand what your private health insurance covers, including exclusions and gaps
    \item Ask your doctor about generic and cheaper medicines
    \item Use the PBS and Medicare safety nets by staying registered and claiming
    \item Know whether you qualify for concession cards and use them
    \item Consider community health and public dental services for lower-cost care
    \item Seek financial counselling if costs are causing stress
    \item Never delay emergency care because of cost
\end{itemize}

\section*{Sources and Further Reading}
The following reputable sources provide further information for healthcare students and patients seeking reliable, current advice about healthcare costs.

\begin{itemize}
    \item Services Australia. \textit{Medicare and the cost of health care.} https://www.servicesaustralia.gov.au/
    \item Australian Government Department of Health and Aged Care. \textit{Pharmaceutical Benefits Scheme.} https://www.pbs.gov.au/
    \item Healthdirect Australia. \textit{Understanding healthcare costs.} https://www.healthdirect.gov.au/
    \item Private Health Insurance Ombudsman. \textit{Understanding private health insurance.} https://www.privatehealth.gov.au/
    \item Consumers Health Forum of Australia. \textit{Health costs and consumer information.} https://chf.org.au/
    \item NHS. \textit{Understanding health costs in the UK.} https://www.nhs.uk/
\end{itemize}
